% !TeX root = ./Outline.tex

\appendix

\section{釋字738號聲請人原因事實整理}

下述整理以\textbf{八件聲請案}為單位；本案共有\textbf{六位自然人提出八件聲請}，八件聲請案分別為陳美英、游輝照、李敬賢（蕃薯藤）、李達緯（金滿意）、李敬賢（頑皮豹）、李達緯（萬華）、吳後田、林蓉。其中李敬賢、李達緯各有兩案。

% 本附錄依下列官方材料整理：釋字第738號解釋本文、司法院於卷內所附之八件聲請書 OCR，
% 以及「釋字第738號解釋其他公開之卷內文書」。
% 下述整理以\textbf{八件聲請案}為單位；本案共有\textbf{六位自然人提出八件聲請}，其中李敬賢、李達緯各有兩案。

% \subsection{聲請案基本結構}

% \begin{itemize}
%   \item 八件聲請案分別為陳美英、游輝照、李敬賢（蕃薯藤）、李達緯（金滿意）、李敬賢（頑皮豹）、李達緯（萬華）、吳後田、林蓉。
%   \item 游輝照案為\textbf{臺北市灣灣電子遊戲場案}，裁判為最高行政法院102年度判字第740號。
%   \item 李敬賢「蕃薯藤」案爭執的是\textbf{臺北縣990公尺距離限制}及其合憲性，裁判為最高行政法院100年度判字第820號。
%   \item 本件多數聲請案的核心，在於\textbf{申請營業級別證、變更登記、或於程序中遭新法加嚴標準阻卻}。
% \end{itemize}

\subsection{八件聲請案總表}

\begin{table}[H]
\centering
\footnotesize
\setlength{\tabcolsep}{3pt}
\caption{釋字738號八件聲請案之基本對照表}
\begin{tabularx}{0.96\textwidth}{L{1.2cm}L{2.0cm}L{1.7cm}Y L{2.4cm}}
\toprule
聲請人 & \celltop{店名 / \\ 地區} & 申請或爭執類型 & 原因事實要點 & 前審與\textbf{確定終局裁判} \\
\midrule
陳美英 & \celltop{金營 / \\新北市三重區 \\（臺北縣三重市）} & 營業場所面積變更 & 原已領有限制級營業級別證，原核准面積77.42平方公尺；99年7月5日申請變更為171.17平方公尺，因新增部分周遭990公尺內有正義國小等9校，依民國95年06月28日公布施行而99年12月25日公告繼續適用之《臺北縣電子遊戲場業設置自治條例》第4條第1項遭否准。 & 北高行99年度訴字第2377號、\textbf{最高行100年度裁字第1601號} \\
\midrule
游輝照 & \celltop{灣灣 / \\臺北市萬華區} & 限制級營業級別證申請 & 99年8月31日先獲准商業設立登記，同日申請限制級營業級別證；100年9月14日第一次否准，101年2月1日訴願撤銷後，101年7月3日改依100年11月10日公布施行之《臺北市電子遊戲場業設置管理自治條例》第5條第1項第2款重為審查並再次否准。 & 北高行102年度訴字第56號、\textbf{最高行102年度判字第740號} \\
\midrule
李敬賢 & \celltop{蕃薯藤 / \\臺北縣轄內} & 距離限制合憲性爭議 & 依所附聲請書，其爭點為臺北縣自治條例第4條之990公尺限制及作業要點援引自治條例之合法性；原處分係依《臺北縣電子遊戲場業設置自治條例》第4條第1項及作業要點審查而否准。 & 北高行98年度訴字第2050號、\textbf{最高行100年度判字第820號} \\
\midrule
李達緯 & \celltop{金滿意 / \\臺北市萬華區} & 限制級營業級別證申請 & 100年6月14日先獲准商業設立登記；100年11月9日申請營業級別證；101年1月20日第一次否准，101年5月訴願撤銷，101年7月3日改依100年11月10日公布施行之《臺北市電子遊戲場業設置管理自治條例》第5條第1項第2款再否准。資本額20萬元，面積116.25平方公尺，亦不符該條例所要求之30公尺道路及1000公尺距離。 & 北高行102年度訴字第106號、\textbf{最高行102年度判字第606號} \\
\midrule
李敬賢 & \celltop{頑皮豹 / \\臺北市萬華區} & 限制級營業級別證申請 & 卷內所附聲請書對應之裁判為北高行102年度訴更一字第145號、最高行103年度裁字第1099號；其否准所依者，為100年11月10日公布施行之《臺北市電子遊戲場業設置管理自治條例》第5條第1項第2款。 & 北高行102年度訴更一字第145號、\textbf{最高行103年度裁字第1099號} \\
\midrule
李達緯 & \celltop{萬華 / \\臺北市萬華區} & 限制級營業級別證申請 & 卷內所附聲請書對應之裁判為北高行102年度訴更一字第135號、最高行104年度判字第35號；其否准所依者，為100年11月10日公布施行之《臺北市電子遊戲場業設置管理自治條例》第5條第1項第2款，仍屬臺北市限制級設立許可受1000公尺、資本額與面積門檻拘束之案件。 & 北高行102年度訴更一字第135號、\textbf{最高行104年度判字第35號} \\
\midrule
吳後田 & \celltop{凱樂喜 / \\新北市三重區 \\（臺北縣三重市）} & 營業場所面積變更 & 聲請書明載：原已營業，申請營業場所面積變更；99年7月21日因周遭990公尺內有學校，依《臺北縣電子遊戲場業設置自治條例》第4條第1項遭否准；爭點與陳美英案極近。 & 北高行99年度訴字第2530號、\textbf{最高行100年度判字第2089號} \\
\midrule
林蓉 & \celltop{廣吉祥 / \\桃園縣} & 級別變更 & 原為既存電子遊戲場，後申請營業級別變更；卷內聲請書及裁判顯示，其否准所依者為民國96年01月04日公布施行之《桃園縣電子遊戲場業設置自治條例》第4條第1項，爭點在於800公尺限制及變更登記時是否應重為審查。 & 北高行100年度訴字第1923號、\textbf{最高行101年度裁字第1794號} \\
\bottomrule
\end{tabularx}
\end{table}

% \subsection{逐案原因事實摘要}

% \subsubsection{陳美英即金營電子遊戲場業}

% % 陳美英案是本件最適合作為母案的事實樣態。

% 依聲請書，
% 陳美英原已在新北市三重區合法經營限制級電子遊戲場，並持有編號09900021之營業級別證，原核准營業面積為77.42平方公尺。嗣於99年7月5日，檢具申請書與建築物使用項目更動執照，申請將營業面積擴大至171.17平方公尺。

% 新北市政府審查後認為，新增部分之營業場所周遭990公尺範圍內有正義國小等9間學校，違反臺北縣電子遊戲場業設置自治條例第4條，遂以99年7月21日北府經登字第0990622195號函否准。其後訴願遭駁回，行政訴訟亦敗訴確定。此案最重要的事實，不在於「新設」本身，而在於法院將\textbf{既有業者擴張營業面積}視同新設部分，因而適用新北（當時臺北縣）990公尺距離限制。

% \subsubsection{游輝照即灣灣電子遊戲場業}

% 游輝照案為\textbf{臺北市新設限制級申請案}。聲請書首頁即載明，游輝照於99年8月31日先經臺北市政府核准商業設立登記，同日申請核發限制級電子遊戲場業營業級別證。臺北市政府先於99年9月17日通知補正，其後於100年9月14日以府產業商字第10032656100號函否准。聲請人提起訴願後，經濟部以101年2月1日經訴字第10106100680號決定撤銷原處分，命另為適法處分。

% 其後，100年11月10日公布施行之《臺北市電子遊戲場業設置管理自治條例》已開始適用。臺北市政府遂改以資本額、道路寬度、1000公尺距離、單層營業面積等門檻重行審查，並於101年7月3日以府產業商字第10133957700號函再次否准。此案的重要性，在於呈現\textbf{申請程序進行中，臺北市之審查基礎由自治條例施行前之規範架構，轉為適用100年11月10日公布施行之自治條例}的處理方式。

% \subsubsection{李敬賢即蕃薯藤電子遊戲場}

% 依本件卷內所附聲請書，蕃薯藤案是以\textbf{臺北縣990公尺距離限制}為爭點的行政訴訟案件。聲請書首頁載明其所不服者，為臺北高等行政法院98年度訴字第2050號判決及最高行政法院100年度判字第820號判決，且其直接挑戰的法規就是臺北縣電子遊戲場業設置自治條例第4條。

% 此案的核心，在於聲請人主張中央法律既已定50公尺標準，地方再以自治條例一律提高至990公尺，已超過管理而構成實質禁絕。

% \subsubsection{李達緯即金滿意電子遊戲場業}

% 金滿意案之聲請人李達緯先於100年6月14日取得商業設立登記，臺北市商業處核准函號為北市商一字第1000006753號。其後於100年11月9日申請電子遊戲場業營業級別證。臺北市政府先於101年1月20日以有礙公共利益等理由否准，經訴願撤銷後，改於101年7月3日以府產業商字第10133965800號函，依100年11月10日公布施行之《臺北市電子遊戲場業設置管理自治條例》重為否准。

% 法院於裁判中認定：金滿意案之商業登記資本額僅20萬元，未達一億元；營業場所臨接道路僅12.73公尺，未達限制級所需之30公尺；營業場所1,000公尺範圍內有多所學校、醫院、圖書館；且申請面積僅116.25平方公尺，未達單層1,000平方公尺。此案清楚顯示臺北市係透過\textbf{資本額、區位、道路、面積}的複合門檻形成整體阻卻效果。

% \subsubsection{李敬賢即頑皮豹電子遊戲場業}

% 頑皮豹案在卷內所附聲請書中，對應的是臺北高等行政法院102年度訴更一字第145號及最高行政法院103年度裁字第1099號。就本件附錄而言，可確定者有二：第一，它與灣灣、金滿意、萬華案同屬\textbf{100年11月10日公布施行之《臺北市電子遊戲場業設置管理自治條例》適用後之限制級申請爭議}；第二，裁判路徑顯示其亦經歷訴願、重為處分及訴訟後，最終仍因該自治條例所設複合門檻而未獲許可。

% 若後續正文需要更細緻的申請地址、資本額與周邊學校清單，宜再直接引頑皮豹聲請書及其附件；但就導讀附錄而言，可先將其與金滿意、萬華案並列為「臺北市複合門檻型案件」。

% \subsubsection{李達緯即萬華電子遊戲場}

% 萬華案與金滿意案相同，均為李達緯所提，且均針對100年11月10日公布施行之《臺北市電子遊戲場業設置管理自治條例》適用後之限制級申請。卷內首頁可直接確認，萬華案之裁判為臺北高等行政法院102年度訴更一字第135號及最高行政法院104年度判字第35號。

% % 其意義在於顯示：即使同一自然人於同一城市提出不同申請，仍可能分別形成獨立之裁判，並被併入釋字738號審理。釋字738號因此是多名業者在不同程序位置、不同申請態樣下共同累積出的違憲審查母體。

% \subsubsection{吳後田即凱樂喜電子遊戲場業}

% 吳後田案與陳美英案最為接近。卷內聲請書首頁即載，其裁判為臺北高等行政法院99年度訴字第2530號判決及最高行政法院100年度判字第2089號判決。從聲請書後段可知，其主張之一亦是：申請變更登記時，除營業場所地址及面積外，本不應一概溯及適用作業要點及自治條例。

% 此外，卷內文字顯示，吳後田主張新北之990公尺限制已近於「實質禁絕」，並批評地方未區分普通級與限制級，一律適用990公尺。就原因事實言，此案屬於\textbf{既有業者擴張面積受阻案}。

% \subsubsection{林蓉即廣吉祥電子遊戲場業}

% 林蓉案對桃園部分最具代表性。卷內聲請書首頁清楚記載，其不服之裁判為臺北高等行政法院100年度訴字第1923號判決與最高行政法院101年度裁字第1794號裁定；聲請內容則直指桃園縣電子遊戲場業設置自治條例第4條第1項之800公尺限制。

% 卷內材料並顯示，林蓉案的重點在於\textbf{既有業者於變更營業級別時，是否應重新受距離限制審查}。廣吉祥案因此屬於級別變更或變更登記爭議，藉此引出桃園縣800公尺距離限制之合憲性。

\subsection{類型整理}

\begin{itemize}
  \item 若就程序型態整理，八件聲請案大致可分為四類：
  既有業者面積變更案（陳美英、吳後田）、
  既有業者級別變更案（林蓉）、
  臺北市限制級新設或新申請案（游輝照、金滿意、頑皮豹、萬華；多與100年11月10日公布施行之《臺北市電子遊戲場業設置管理自治條例》之適用有關），
  以及以臺北縣990公尺限制為核心的蕃薯藤案。
  \item 若就地方別整理，臺北市案件呈現的是\textbf{複合門檻化}；且在研究上宜區分\textbf{100年11月10日自治條例公布施行前之審查基礎}與\textbf{該自治條例施行後之複合門檻}。除1000公尺外，後者還同時要求資本額、道路寬度與樓地板面積；新北市案件則更集中於\textbf{990公尺距離限制對變更登記的阻卻效果}；桃園縣案件則集中於\textbf{800公尺距離限制對變更或級別變更的阻卻效果}。
  % \item 
\end{itemize}

\subsection{實證補充}


釋字第738號的事實底層，是\textbf{地方自治條例與作業要點共同作用，讓既有業者之擴張、級別變更，及新申請人之進入市場，均受到強力阻卻。}

\begin{table}[H]
\centering
\footnotesize
\setlength{\tabcolsep}{3pt}
\caption{卷內文書所附 98--103 年地方電子遊戲場業統計摘錄}
\begin{tabularx}{0.96\textwidth}{L{1cm}L{2cm}L{2cm}L{2.1cm}L{2.2cm}Y}
\toprule
年度 & 新北市合法總家數（限） & 新北市總申請件數（限） & 新北市因距離否准件數（限） & 桃園市合法總家數（普/限） & 桃園市因距離否准件數（普/限） \\
\midrule
98年 & 39 & 3 & 0 & 5 / 83 & 0 / 0 \\
\midrule
99年 & 38 & 1 & 0 & 5 / 79 & 0 / 0 \\
\midrule
100年 & 38 & 0 & 0 & 5 / 71 & 0 / 1 \\
\midrule
101年 & 34 & 6 & 0 & 4 / 65 & 0 / 0 \\
\midrule
102年 & 32 & 36 & 36 & 4 / 58 & 0 / 0 \\
\midrule
103年 & 28 & 4 & 4 & 4 / 51 & 0 / 0 \\
\bottomrule
\end{tabularx}
\end{table}

上表可直接用來支持一個重要觀察：就卷內文書所附統計而言，新北市在102年、103年之否准，幾乎全部集中於距離限制；這也是大法官於解釋文末特別提醒地方應隨時檢討，以免產生\textbf{實質阻絕效果}的背景。

\section{釋字738號所涉中央與地方管制規範對照表}

下表依釋字第738號解釋文、理由書及該案所涉及法規文字整理，目的在於呈現本案中中央法律、中央行政命令與地方自治條例之規範層級、規範對象與距離限制之具體差異。



\subsection{本案所涉條文原文摘錄}

下表摘錄本案直接審查或在審查過程中被援用之規範原文。

\begin{center}
\footnotesize
\setlength{\tabcolsep}{4pt}
\captionof{table}{釋字738號所涉條文原文摘錄}
\begin{tabularx}{0.94\linewidth}{L{3.6cm}Y L{3.2cm}}
\toprule
規範 & 原文摘錄 & 本案中的作用 \\
\midrule
電子遊戲場業管理條例第9條第1項 &
「電子遊戲場業之營業場所，應距離國民中、小學、高中、職校、醫院五十公尺以上。」 &
中央法律所設最低距離標準；大法官據此判斷地方是否仍有加嚴空間。 \\
\midrule

電子遊戲場業申請核發電子遊戲場業營業級別證作業要點第2點第1款第1目 &
「申請作業程序：電子遊戲場業……申請電子遊戲場業營業級別證或變更登記，應符合下列規定：（一）營業場所：1.符合電子遊戲場業管理條例、自治條例及其他有關規定。」 &
中央行政規則；將自治條例納入級別證申請及變更登記之審查架構。 \\
\midrule

臺北市電子遊戲場業設置管理自治條例第5條第1項第2款 &
「電子遊戲場業之營業場所應符合下列規定：……二、限制級：……應距離幼稚園、國民中、小學、高中、職校、醫院、圖書館一千公尺以上。」 &
臺北市限制級電子遊戲場之距離管制依據；並增加幼稚園、圖書館為保護場所。 \\
\midrule

臺北縣電子遊戲場業設置自治條例第4條第1項 &
「前條營業場所，應距離國民中、小學、高中、職校、醫院九百九十公尺以上。」 &
臺北縣原規範，改制後由新北市繼續適用；對普通級、限制級均形成990公尺限制。 \\
\midrule

桃園縣電子遊戲場業設置自治條例第4條第1項 &
「電子遊戲場業之營業場所，應距離國民中、小學、高中、職校、醫院八百公尺以上。」 &
桃園縣距離限制之直接法源；在林蓉案中用以審查級別變更申請。 \\
\bottomrule
\end{tabularx}
\end{center}




\subsection{規範差異摘要}

\begin{itemize}
  \item 中央法律本身僅明定電子遊戲場業應與國民中、小學、高中、職校及醫院保持50公尺以上距離，且未區分普通級與限制級。
  \item 中央作業要點本身不另外設定公尺數，而是將申請核發營業級別證或變更登記時應符合自治條例等規範納入審查架構。
  \item 臺北市自治條例僅就\textbf{限制級}明定1,000公尺距離，並將幼稚園、圖書館納入保護場所。
  \item 臺北縣自治條例對\textbf{普通級與限制級一體適用}990公尺限制，保護場所為國民中、小學、高中、職校及醫院。
  \item 桃園縣自治條例亦未區分普通級與限制級，採800公尺限制，保護場所為國民中、小學、高中、職校及醫院。
\end{itemize}




\begin{landscape}
\thispagestyle{plain}
\begin{center}
\footnotesize
\setlength{\tabcolsep}{4pt}
\captionof{table}{釋字738號所涉中央與地方相關管制規範對照表}
\begin{tabularx}{0.94\linewidth}{L{1.6cm}L{4.6cm}Y L{2.1cm}L{1.8cm}L{5.3cm}}
\toprule
層級 & \celltop{條例名稱 /\\條號} & 場所類型與距離限制 & \celltop{普通級 /\\限制級} & 法律性質 & 備註 \\
\midrule
中央 &
\celltop{電子遊戲場業管理條例\\第9條第1項} &
電子遊戲場業之營業場所，應距離國民中、小學、高中、職校、醫院50公尺以上。 &
未區分 &
法律 &
解釋理由書將其定位為法律保留地方因地制宜空間所設之最低標準。 \\
\midrule

中央 &
\celltop{電子遊戲場業申請核發電子遊戲場業營業級別證作業要點\\第2點第1款第1目} &
申請營業級別證或變更登記時，營業場所應符合電子遊戲場業管理條例、自治條例及其他有關規定；本身未直接規定固定公尺數。 &
未區分 &
行政規則 &
解釋理由書認其僅係指明申請時應適用之法令，屬細節性、技術性規定。 \\
\midrule

臺北市 &
\celltop{臺北市電子遊戲場業設置管理自治條例\\第5條第1項第2款} &
限制級電子遊戲場業應距離幼稚園、國民中、小學、高中、職校、醫院、圖書館1,000公尺以上。 &
限制級 &
自治條例 &
除距離加嚴外，規範對象另擴及幼稚園、圖書館。 \\
\midrule

\celltop{臺北縣／\\新北市} &
\celltop{臺北縣電子遊戲場業設置自治條例\\第4條第1項} &
前條營業場所應距離國民中、小學、高中、職校、醫院990公尺以上。 &
普通級與限制級均包括在內 &
自治條例 &
原為臺北縣自治條例；改制後由新北市繼續適用，後因期限屆滿失效。 \\
\midrule

桃園縣 &
\celltop{桃園縣電子遊戲場業設置自治條例\\第4條第1項} &
電子遊戲場業之營業場所，應距離國民中、小學、高中、職校、醫院800公尺以上。 &
未區分 &
自治條例 &
- \\
\bottomrule
\end{tabularx}
\end{center}
% \end{landscape}

% \begin{landscape}
% \thispagestyle{plain}
\begin{center}
\footnotesize
\setlength{\tabcolsep}{4pt}
\captionof{table}{釋字738號中央與地方距離管制之摘要比較}
\begin{tabularx}{0.92\linewidth}{L{2.8cm}L{4.8cm}L{3.6cm}L{2.2cm}L{3.2cm}}
\toprule
\celltop{中央或地方} & 共同場所類型 & 另增場所類型 & 距離 & \celltop{普通級 /\\限制級} \\
\midrule
\celltop{中央（法律）} & 國民中、小學、高中、職校、醫院 & 無 & 50公尺以上 & 未區分 \\
\midrule
% \makecell[l]{中央\\（作業要點）} & 申請時應符合中央法律、自治條例及其他有關規定 & 無獨立增列 & \makecell[l]{未自定\\固定公尺數} & 未區分 \\
臺北市 & 國民中、小學、高中、職校、醫院 & 幼稚園、圖書館 & 1,000公尺以上 & 僅限制級明定距離限制 \\
\midrule
臺北縣／新北市 & 國民中、小學、高中、職校、醫院 & 無 & 990公尺以上 & 普通級、限制級均適用 \\
\midrule
桃園縣 & 國民中、小學、高中、職校、醫院 & 無 & 800公尺以上 & 未區分 \\
\bottomrule
\end{tabularx}
\end{center}
\end{landscape}



% \begin{landscape}
% \thispagestyle{plain}


% \end{landscape}
